% Vietnamese translation for OI Public Library
% Source-File: IOI中国国家候选队论文/2008/Day2/6.余林韵《运用化归思想解决信息学中的数列问题》/余林韵_运用化归思想解决信息学中的数列问题.ppt
\documentclass[11pt]{article}
\usepackage{oipl-vn}
\usepackage{booktabs}

\newcommand{\Oh}{\mathcal{O}}

\title{Bản trình chiếu: Quy chuyển trong các bài toán dãy số}
\author{Yu Linyun}
\date{2008}

\begin{document}
\maketitle

\origtitle{Reduction in sequence problems}
\sourcepath{IOI China candidate team papers / 2008 / Day 2 / Yu Linyun.ppt}

\section{Phạm vi bản dịch}

Bản PowerPoint đi kèm bài viết và bản thuyết trình về việc dùng tư tưởng quy
chuyển để giải các bài toán dãy số. Phần chữ trích được từ PPT lặp lại các
mốc chính: công thức của \textit{Dispute}, modulo \(9973\), TopCoder SRM 369,
\textit{Count} của NOI 2007, \textit{Maximum. Version 2} và Ural 1396. Bản
ghi chú này đi theo các mốc slide đó, rút gọn những chứng minh dài trong bài
viết chính nhưng giữ lại ý tưởng thuật toán.

Thông điệp của slide là: với dãy số, không nên chỉ sinh vài số hạng đầu rồi
đoán quy luật. Cần xác định đối tượng quy chuyển, mục tiêu quy chuyển và con
đường quy chuyển; các tính chất như tuần hoàn, tuyến tính theo từng đoạn,
trạng thái hữu hạn hoặc cấu trúc Fibonacci thường là chìa khóa.

\section{Tư tưởng quy chuyển}

Quy chuyển là biến một bài toán mới thành một bài toán đã biết cách giải hoặc
dễ xử lý hơn. Trong ngữ cảnh dãy số, các đích quy chuyển thường gặp là:
\begin{itemize}
  \item đưa truy hồi về truy hồi tuyến tính hệ số hằng để dùng nhân ma trận;
  \item tìm chu kỳ modulo hoặc chu kỳ cuối cùng của dãy;
  \item chia quá trình thành các giai đoạn giống thuật toán Euclid;
  \item thay một dãy đơn bằng một vector trạng thái rồi lũy thừa ma trận
  chuyển;
  \item nhận ra cấu trúc Fibonacci hoặc cấp số trong các đoạn đặc biệt.
\end{itemize}

Các slide vì vậy không trình bày dãy số như một chủ đề thuần toán học, mà xem
mỗi dãy là một mô hình cần được đổi dạng cho phù hợp với ràng buộc dữ liệu.

\section{\textit{Dispute}: từ truy hồi phi tuyến đến bước nhảy tuyến tính}

Ví dụ mở đầu là dãy
\[
  F_0=0,\qquad F_n=g(n,F_{n-1}),
\]
trong đó slide nhấn mạnh công thức
\[
  g(x,y)=\bigl((y-1)x^5+x^3-xy+3x+7y\bigr)\bmod 9973.
\]
Nếu chỉ nhìn trực tiếp, \(n\) lớn làm truy hồi từng bước không thể chấp nhận.
Một cách thô là lưu các mốc \(F_{100000k}\), rồi từ mốc gần nhất truy hồi
thêm không quá \(100000\) bước. Ý tưởng ấy đúng hướng ở chỗ nó giảm quy mô,
nhưng chưa dùng cấu trúc của công thức.

Điểm đột phá nằm ở hai quan sát mà slide tách riêng:
\begin{enumerate}
  \item với \(x\) cố định, công thức tuyến tính theo \(y\);
  \item mọi tính toán đều lấy modulo số nguyên tố \(M=9973\).
\end{enumerate}
Do \(g(x,y)\) chỉ phụ thuộc vào \(x\bmod M\), khi nhảy đúng \(M\) bước ta thu
được hai hằng số \(A,B\) sao cho
\[
  F_{(k+1)M}\equiv AF_{kM}+B\pmod M.
\]
Như vậy dãy con ở các chỉ số bội của \(M\) đã trở thành truy hồi tuyến tính
hệ số hằng. Ta tính \(A,B\) trong \(\Oh(M)\), dùng lũy thừa nhanh hoặc công
thức tổng cấp số nhân modulo để lấy \(F_{\lfloor n/M\rfloor M}\), rồi truy
hồi thêm phần dư. Độ phức tạp được rút còn
\[
  \Oh\!\left(M+\log_2\frac{n}{M}\right).
\]

\section{TopCoder SRM 369: chu kỳ và thuật toán Euclid}

Mốc \textit{TopCoder SRM 369} trong slide tương ứng với ví dụ dãy trị tuyệt
đối:
\[
  F_1=A,\qquad F_2=B,\qquad F_n=|F_{n-1}-F_{n-2}|.
\]
Miền \(n\) rất lớn nên không thể sinh dãy. Quan sát nhiều cặp ban đầu cho
thấy cuối cùng dãy rơi vào chu kỳ ba phần tử
\[
  x,\ x,\ 0,\ x,\ x,\ 0,\ldots
\]
và \(x=\gcd(A,B)\).

Chứng minh trong bài viết chính đi theo đúng tinh thần quy chuyển: mọi số hạng
đều là bội của \(\gcd(A,B)\), giá trị cực đại của hai số liên tiếp giảm sau
một số bước, và khi hai số liên tiếp bằng nhau thì chu kỳ cuối xuất hiện. Từ
đó quá trình được chia thành các giai đoạn giống thuật toán Euclid. Mỗi giai
đoạn bắt đầu bằng \(x>y\), đặt \(\Delta=x-y\), rồi nhảy qua cả một cụm số hạng
thay vì trừ từng bước. Thời gian còn \(\Oh(\log n)\).

\section{\textit{Count} của NOI 2007: một dãy thành vector trạng thái}

Slide tiếp theo nêu \textit{Count} của NOI 2007. Với \(n\) đỉnh trên một đường
thẳng, nối hai đỉnh nếu khoảng cách không quá \(k\); cần đếm số cây khung
modulo \(65521\). Khi \(k=2\), số nghiệm là các số Fibonacci ở vị trí chẵn,
nhưng với \(k=3,4,5\) thì dãy đơn không còn dễ nhận dạng.

Quy chuyển ở đây là đổi đối tượng: thay vì chỉ giữ số cây khung của \(n\) đỉnh,
ta giữ trạng thái liên thông của \(k\) đỉnh cuối. Hai vị trí cùng số nghĩa là
thuộc cùng một thành phần. Ví dụ với \(k=3\), các trạng thái chuẩn là
\[
  111,\quad 112,\quad 121,\quad 122,\quad 123.
\]
Khi thêm đỉnh mới, ta thử các cạnh nối tới \(k\) đỉnh trước đó, loại những
cạnh tạo chu trình, rồi chuẩn hóa lại trạng thái. Vì \(k\le5\), số trạng thái
lớn nhất chỉ là \(52\). Ta xây dựng ma trận chuyển \(52\times52\), lấy lũy
thừa \(n-k\), và đọc trạng thái mà \(k\) đỉnh cuối cùng liên thông.

Độ phức tạp slide nhấn mạnh là
\[
  \Oh(52^3\log n)
\]
với bộ nhớ \(\Oh(52^2)\). Đây là ví dụ quan trọng vì nó cho thấy ``dãy số''
trong đề bài có thể ẩn một hệ nhiều dãy liên kết với nhau.

\section{\textit{Maximum}: cấu trúc Fibonacci trong từng đoạn}

Mốc cuối của PPT là \textit{Maximum. Version 2} / Ural 1396. Dãy được cho bởi
\[
  A_1=1,\qquad A_{2i}=A_i,\qquad A_{2i+1}=A_i+A_{i+1}.
\]
Với \(n\), cần tìm giá trị lớn nhất trong \(A_1,\ldots,A_n\).

Vì công thức gắn với phép nhân đôi chỉ số, slide chia dãy theo các đoạn giữa
hai lũy thừa của \(2\). Các đoạn đầu là:
\[
\begin{array}{l}
1\\
1,\ 2\\
1,\ 3,\ 2,\ 3\\
1,\ 4,\ 3,\ 5,\ 2,\ 5,\ 3,\ 4.
\end{array}
\]
Giá trị lớn nhất của mỗi đoạn lần lượt là Fibonacci. Tổng quát hơn, nếu trong
một đoạn có hai số kề \(a,b\), thì sau khi mở rộng thêm \(t\) tầng, giá trị
lớn nhất trong phần sinh ra là một số hạng của dãy Fibonacci tổng quát bắt
đầu từ \(a,b\). Nhờ đó, lời giải chỉ cần đi theo biểu diễn nhị phân của \(n\)
và cập nhật cặp biên, đạt \(\Oh(\log n)\).

\section{Kết luận}

Bốn ví dụ trong slide dùng bốn con đường quy chuyển khác nhau: bước nhảy
modulo để tuyến tính hóa, giai đoạn Euclid, vector trạng thái và cấu trúc
Fibonacci theo tầng. Điểm chung là không cố xử lý dãy ở dạng ban đầu. Ta tìm
một hình thức biểu diễn làm lộ ra quan hệ đã biết, rồi dùng công cụ tiêu chuẩn
như Euclid, nhân ma trận hoặc lũy thừa nhanh để xử lý miền dữ liệu lớn.

\end{document}
